% TOP_PROBLEM: {"id": 20000584, "title": "Higher-dimensional quantitative bi-Lipschitz parameterization and a scalar-curvature obstruction", "category_id": 8, "category": {"id": 8, "name": "analysis", "display_name": "Analysis", "description": "Limits, continuity, calculus, and function theory.", "slug": "analysis", "order_index": 8, "created_at": "2026-07-31T15:26:25.671Z"}, "status": "partially_solved"}
% FIRST_POSTED: 2026-10-06
% ATTEMPT_STATUS: unresolved
% OPENAI_CLAIM: {"schema_version":1,"source_repo":"https://github.com/openai/math","source_commit":"adc7f1241b42e322a6451854ab7e4b4c146bf78a","release_date":"2026-10-06","match":"partial","resolution":"partial","independently_reviewed":false,"second_pass":"n/a","families":[{"family":"357","title":"Bi-Lipschitz coordinates at every regular RCD point","subject":"Differential geometry","manuscripts":[{"title":"Bi-Lipschitz Coordinates at Regular Points of Noncollapsed RCD Spaces","pdf_path":"preprints/Bi-Lipschitz-Coordinates-at-Regular-Points-of-Noncollapsed-RCD-Spaces-September-25-2026/paper.pdf","readme_path":"preprints/Bi-Lipschitz-Coordinates-at-Regular-Points-of-Noncollapsed-RCD-Spaces-September-25-2026/README.md","date":"2026-09-25","theorem_ref":"build/source/introduction.tex:17, Theorem thm:main"}],"lean":null,"reasoning_trace":null}]}
\documentclass[11pt]{article}
\usepackage[margin=25mm]{geometry}
\usepackage{fontspec}
\setmainfont{Latin Modern Roman}
\usepackage{amsmath,amssymb,mathtools}
\usepackage{unicode-math}
\setmathfont{Latin Modern Math}
\usepackage{enumitem}
\usepackage{xurl,hyperref,bookmark}
\hypersetup{colorlinks=true,urlcolor=blue}
\setcounter{secnumdepth}{0}
\setlength{\parindent}{0pt}
\setlength{\parskip}{5pt}
\setlength{\emergencystretch}{3em}
\begin{document}
\section*{Higher-dimensional quantitative bi-Lipschitz parameterization and a scalar-curvature obstruction}

\subsection{OpenAI's claim}
\textit{Bi-Lipschitz charts at regular RCD points}: For noncollapsed $\mathrm{RCD}(K,n)$ spaces with $n\ge2$, every regular point is claimed to have an open Euclidean bi-Lipschitz neighborhood. The distortion bound depends only on dimension, and the neighborhood carries the restricted ambient distance.

\subsection{Outline of the argument}
\textit{Bi-Lipschitz charts at regular RCD points}: Poisson replacements and two-center estimates produce summable directional defects for local metric matrices. An iteration on a common open coordinate image corrects the pullback metric using inverse-transpose Jacobians and smooth covector fields. Diagonal second-derivative recurrences and higher-derivative bounds keep the corrections summable even when their square roots are not, yielding invertible bi-Lipschitz coordinates.

\subsection{Scope relative to this problem}
\textit{Bi-Lipschitz charts at regular RCD points}: The theorem gives quantitative bi-Lipschitz coordinates in every dimension at regular points of noncollapsed RCD spaces. The source asks a wider higher-dimensional Bonk-Lang analogue, including global parameterization, and does not limit itself to such local regular RCD points.

This page records partial support. Complete resolution of the catalogue question is not established by this record.

\subsection{Formal verification}
\textit{Bi-Lipschitz charts at regular RCD points}: The pinned catalogue lists no Lean formalization for this result.

This site has not run the Lean development. These links describe OpenAI’s selected formal statements and should be read with their stated scope.

\subsection{Status caveat}
OpenAI reports that an internal model produced the manuscripts. Their arguments have not been independently verified as proofs by this site. Any portion outside the documented formal scope may contain mistakes, and the catalogue label reports the announced mathematical claim.

\subsection{Sources}
\begin{itemize}
\item OpenAI, \emph{Bi-Lipschitz Coordinates at Regular Points of Noncollapsed RCD Spaces}, 2026-09-25; citation key \texttt{OAI:Bi-Lipschitz-Coordinates-at-Regular-Points-of-Noncollapsed-RCD-Spaces-September-25-2026}. \href{https://github.com/openai/math/blob/adc7f1241b42e322a6451854ab7e4b4c146bf78a/preprints/Bi-Lipschitz-Coordinates-at-Regular-Points-of-Noncollapsed-RCD-Spaces-September-25-2026/paper.pdf}{PDF (GitHub)}; \href{https://github.com/openai/math/raw/adc7f1241b42e322a6451854ab7e4b4c146bf78a/preprints/Bi-Lipschitz-Coordinates-at-Regular-Points-of-Noncollapsed-RCD-Spaces-September-25-2026/paper.pdf}{PDF (direct)}; \href{https://github.com/openai/math/blob/adc7f1241b42e322a6451854ab7e4b4c146bf78a/preprints/Bi-Lipschitz-Coordinates-at-Regular-Points-of-Noncollapsed-RCD-Spaces-September-25-2026/README.md}{Manuscript page and citation}; \href{https://github.com/openai/math/blob/main/preprints/Bi-Lipschitz-Coordinates-at-Regular-Points-of-Noncollapsed-RCD-Spaces-September-25-2026/paper.pdf}{Latest version}.
\item \href{https://github.com/openai/math/blob/adc7f1241b42e322a6451854ab7e4b4c146bf78a/CONTENTS.md}{OpenAI catalogue entry}.
\end{itemize}
\end{document}
